\section{Introduction}

\textit{\textbf{Address the motivation to the project.} What the reader should know before going into more details.}

\begin{itemize}
  \item \textbf{Event cameras have moved from a research curiosity to commercial deployment in robotics, automotive, and AR/VR, driven by demands for low latency and high dynamic range (Gallego et al., 2022).} This growing adoption spans several distinct application clusters: SLAM and obstacle avoidance for robotics spatial awareness (Vidal et al., 2018), perception for autonomous driving and driver-assistance systems in automotive (Maqueda et al., 2018), and low-power always-on eye tracking for AR/VR (Angelopoulos et al., 2021), applications where conventional frame-based sensors struggle to keep up with fast motion or extreme lighting.
  \item \textbf{Unlike a frame camera that samples all pixels at a fixed rate, a DVS pixel independently and asynchronously emits an "event" whenever its log-brightness changes by a threshold amount, producing a stream of (x, y, t, polarity) tuples (Lichtsteiner, Posch \& Delbruck, 2008; Gallego et al., 2022).} This per-pixel, change-triggered operating principle is the mechanism behind every property discussed next, and it is also why event data cannot simply be handled with tools built for frames.
  \item \textbf{As a result of this sensor design, event cameras exhibit unique features: microsecond temporal resolution, very high dynamic range (\textasciitilde{}140 dB vs \textasciitilde{}60 dB), low power, and no motion blur; but also an output format that conventional vision algorithms can't consume directly (Gallego et al., 2022).} These properties let event cameras operate reliably in conditions that defeat frame-based sensors, capturing fast motion without blur and remaining usable in scenes with extreme contrast between shadow and highlight, but the same asynchronous, sparse output that enables this performance has no fixed grid or sampling rate, so it cannot be fed directly into vision algorithms built around dense image frames.
  \item \textbf{Because events aren't frames, they are either converted into intermediate representations or fed directly to asynchronous/spiking methods.} Common representation choices include binned event frames (Maqueda et al., 2018), binned binary event frames (Barchid, Mennesson \& Djéraba, 2022), memory surfaces (Afshar et al., 2019), and voxel grids (Zhu et al., 2019), each repackaging the asynchronous stream into a format compatible with standard deep-learning architectures.
  \item \textbf{Most current machine learning methods built on these representations are supervised solutions.} Supervised training needs large labeled datasets, but real event recordings are costly to capture and annotate (high-speed ground truth is difficult to obtain), which is the standard data bottleneck cited across the field (Wang et al., 2024) and echoed by simulator papers such as v2e (Hu et al., 2021) and DVS-Voltmeter (Lin et al., 2022) as core motivation for synthetic data.
  \item \textbf{Collecting real data being expensive pushes the field toward synthetic data generated by sensor simulators, and v2e remains the most widely used despite known limitations (Hu et al., 2021). v2e takes only frame input and relies on SuperSlowMotion interpolation to synthesize intermediate frames, which is inefficient and produces layered events; DVS-Voltmeter (Lin et al., 2022), PECS (Han et al., 2024), and V2CE (Zhang et al., 2024) each address a specific piece of this pipeline.}
  \item \textbf{But there is always a domain gap between real and simulated data, which degrades synthetic data quality, and every one of these simulators, including v2e's successors, is validated qualitatively or against a specific downstream task rather than against a common, direct measure of realism (Tan, Kumar B N \& Chakravarthi, 2025; Chen \& Negrut, 2024).} This means realism claims can't be compared across methods on equal footing; recent learned metrics such as EQS (Chanda et al., 2025) begin to close that gap, but they depend on features from a pretrained network rather than a physically interpretable, reproducible measurement.
  \item \textbf{We propose a metric framework for directly quantifying the real-vs-synthetic event-data gap, combining multi-kernel MMD (Gretton et al., 2012) and Sliced Wasserstein Distance (Bonneel et al., 2015; Flamary et al., 2021), benchmarked against a Chamfer-distance baseline, and validated on recorded optical-chopper data on a Prophesee EVK4.}
  \item \textbf{Road map of the paper.} Briefly preview the section order (background, methodology including the framework overview, evaluation protocol, and the chopper/saccade results) so the reader knows what's coming before the detail begins.
\end{itemize}

\section{Background}

\textit{\textbf{Establish the concepts the reader needs before the method.} If the venue forces a shorter format, fold this into "Related Work + Preliminaries" (see To do).}

\begin{itemize}
  \item \textbf{The ideal DVS pixel model is a brightness change detector.} Each pixel continuously monitors log-brightness; when the change since the last event crosses a contrast threshold θ, it emits an ON or OFF event and enters a brief refractory period. This gives the generative model that both the sensor and any faithful simulator must follow.
  \item \textbf{Real sensors depart from this ideal model in the first place because θ is not a quantity the user sets directly: bias settings are specific to systems, not thresholds that are directly interpretable with physical measurements.} The user-facing bias settings (e.g., bias\_diff\_on/off) are currents that shift circuit operating points, and their mapping to the effective θ depends on factory trim, temperature, and photocurrent, so a bias value cannot be read as a threshold value.
  \item \textbf{Beyond the threshold itself, several further non-idealities keep a physically accurate camera model out of reach.} Pixel-to-pixel threshold mismatch, intensity-dependent bandwidth and latency, and background-activity noise are the dominant departures from the ideal model, and they are exactly the effects a simulator tries to reproduce and a similarity metric must therefore be sensitive to.
  \item \textbf{Reproducing these effects faithfully is the simulator's task; measuring whether it has done so requires a representation of the event stream that is neutral to any particular application, and events lend themselves naturally to being treated as spatiotemporal point sets.} A recording is a set of (x, y, t, p) tuples, which can be viewed as samples drawn from an underlying distribution over a 3D space-time volume rather than as a signal to be decoded, and this shift in viewpoint is what allows distribution-distance machinery to apply directly to raw events.
  \item \textbf{Once both real and synthetic recordings are point sets from an underlying distribution, measuring the domain gap becomes a two-sample problem.} Given a real event set X and a synthetic set Y from the same stimulus, "how realistic is Y?" turns into the statistical question "were X and Y drawn from the same distribution?", which can be answered from the event data alone without training or evaluating any downstream task.
  \item \textbf{The first tool for answering that question is Maximum Mean Discrepancy (MMD), which compares distributions through their mean embeddings in a kernel space.} Each distribution is mapped into a reproducing kernel Hilbert space and represented by the mean of its embedded samples; MMD is the distance between those two means, so it captures differences in the full structure of the distributions rather than in any single summary statistic.
  \item \textbf{The second tool, Wasserstein distance, measures the same discrepancy through mass transport rather than kernel embeddings.} Wasserstein distance is the minimal "work" required to move one distribution's mass onto the other's, and its sliced variant (SWD) replaces the expensive high-dimensional transport problem with an average of one-dimensional Wasserstein distances taken over many projection directions, which is what makes it computable at event-stream scale.
  \item \textbf{Against these two distribution-level metrics we set Chamfer distance, a point-level baseline that makes no distributional assumption at all.} Chamfer distance sums the distance from each point in one set to its nearest neighbour in the other, which is intuitive and widely used in point-cloud work, but it compares point positions rather than distributions and so is included here as the baseline our distribution-level metrics are judged against.
  \item \textbf{Spike-train metrics from computational neuroscience, such as the van Rossum distance and SPIKE-distance, are a natural-seeming alternative but do not transfer to this setting.} The van Rossum distance requires convolving every spike train with a kernel and integrating pairwise, which is computationally prohibitive at the event counts a 1280×720 sensor produces, while SPIKE-distance is formulated around the timing structure of one-dimensional trains and does not extend naturally to a stream whose spatial arrangement carries the signal.
\end{itemize}

\section{Methodology}

\subsection{Framework overview}

\begin{itemize}
  \item \textbf{The framework runs as a fixed pipeline from paired acquisition through to a set of distance values per condition.} Real and synthetic events are acquired in pairs from the same stimulus, then converted into a common point-set representation that is polarity-split, space-time scaled, and windowed. Two distribution-level metrics, MMD and SWD, are computed on that representation alongside a Chamfer baseline. The resulting values are validated first on controlled stimuli and then applied in the eye-tracking case study.
  \item \textbf{The first stage, paired acquisition, is what makes the residual difference attributable to the simulator rather than to the scene.} An EVK4 event camera and a Mako frame camera view the same stimulus, both externally triggered from a common source so that the two streams are synchronised, and the Mako frames then drive v2e, so real and synthetic events originate from an identical physical scene and any difference between them is a property of the simulator rather than of what was filmed.
  \item \textbf{For validation the stimulus is an optical chopper, chosen because it produces a periodic signal at a known, tunable frequency.} Each trial contains many nominally identical periods, which supplies repeated samples of the same stimulus for estimating variability and a noise floor, and sweeping the chopper frequency gives a single physically controlled variable along which the metrics can be tested. The chopper motor is free-running rather than phase-locked to the cameras, so periods are recovered from its reference output rather than assumed from frame count.
  \item \textbf{Throughout the pipeline, every free parameter is fixed by rule rather than estimated from the data being compared.} Kernel bandwidths, projection count, and the space-time unit conversion are set from sensor geometry or application physics, not fitted to the two sets under comparison, so a number computed under one condition is comparable to a number computed under another.
  \item \textbf{The measurement is defined for a single sensor at a time: it compares real recordings against a simulator's attempt to reproduce that same sensor.} Within a fixed sensor and configuration the values are directly comparable, which is what lets a simulator be tracked across stimulus conditions or several simulators be ranked against the same recordings; comparing values obtained on different sensors is outside the scope of the framework, since the kernel bandwidths and the space-time scale factor are anchored to the geometry of the sensor being modelled.
\end{itemize}

\subsection{Event representation}

\begin{itemize}
  \item \textbf{To measure distances, we are processing event data from stream to point cloud.} Events in a window become 3D points (x, y, s·t), with s the µs-per-pixel scale factor; ON and OFF events are separated into channels upstream, compared per-channel, then combined.
  \item \textbf{To make the measurement interpretable, meaningful and to avoid having one dimension dominates the result.} We are making space and time commensurable. Any metric on (x, y, t) points needs a µs-per-pixel conversion. We anchor it in the application: the saccade main-sequence relation D = 2.2A + 21 yields \textasciitilde{}21–42 µs/px for our sensor geometry, a physically-motivated rather than data-fitted choice. The amplitude-duration form used here is the relation that remains close to linear over the amplitude range of interest, unlike the amplitude-peak-velocity form which saturates nonlinearly; note also that no single model of the main sequence is universally agreed, which is why the sensitivity check over the 21–42 µs/px range is reported rather than treated as optional.
  \item \textbf{Making data comparison in a fixed temporal window is better than comparing the whole dataset directly.} Window duration and placement policy; note that all downstream parameters are window-independent by construction (motivates the fixed-bandwidth choice below).
\end{itemize}

\subsection{MMD}

\begin{itemize}
  \item \textbf{Maximum Mean Discrepancy (MMD) is a statistical distance metric used to determine if two sets of data are drawn from the same underlying probability distribution.} Instead of comparing individual data points, MMD compares the global shape and structure of the two datasets.
  \item \textbf{Concretely, MMD² is built from three averages of pairwise similarities: how similar the real points are to each other, how similar the synthetic points are to each other, and how similar the two sets are to each other.} If the two sets come from the same distribution, the cross-set similarity is as large as the within-set similarities and the three terms cancel; if they differ, the within-set terms dominate and MMD² grows. We use the unbiased estimator:
  \item \textbf{Every symbol above is defined as follows.} X = \{x₁, …, x\_m\} is the set of real events and Y = \{y₁, …, y\_n\} the set of synthetic events in the window being compared, each point being a 3D coordinate (x, y, s·t) as defined in 3.2; m and n are the number of real and synthetic points, which need not be equal; k(a, b) is the kernel, a similarity function returning a value near 1 when two points are close and near 0 when far apart; and the notation i ≠ j means the sum runs over distinct pairs only, excluding the case where a point is paired with itself. We use the Gaussian kernel k(a, b) = exp(−‖a − b‖² / 2σ²), in which ‖a − b‖ is the Euclidean distance between the two points and σ is the bandwidth, the length scale that sets what counts as "close"; the choice of σ is the subject of the next bullet.
  \item \textbf{The unbiased variant is required here because the biased alternative makes the measurement depend on how many events a recording happens to contain.} The biased estimator divides by m² and n² and includes the self-pairs i = j, for which the Gaussian kernel returns exactly 1; this adds a positive offset that shrinks as the sets grow rather than vanishing, so its value shifts with event count even when the underlying distributions are unchanged. Since event count is not held fixed across our conditions, and in particular rises with sensor noise, that offset would be confounded with the very difference the metric is meant to detect, and in practice produces a spurious U-shaped response as noise inflates the count. Excluding the diagonal removes the offset and makes the estimator unbiased for the true MMD².
  \item \textbf{The implementation has a multi-kernel design with fixed bandwidths.} A small set of bandwidths anchored to sensor geometry (pixel pitch through sensor-scale structure), replacing the median heuristic, which is data- and window-duration-dependent and breaks cross-condition comparability.
  \item \textbf{Small bandwidths are necessary to resolve fine-scale structure, but they carry a failure mode that has to be stated: under added uniform noise the MMD² of a small kernel can decrease rather than increase.} With a small σ the kernel is effectively zero for any pair of points that are not close together, so uniformly distributed noise contributes near-zero similarity to every term while still inflating the pair counts, and the statistic can fall even though the two sets have genuinely become more different. This is the reason the bandwidth set spans several scales rather than reaching for the smallest kernel that resolves the structure of interest, and it is why the response of each kernel to controlled perturbation is characterised before the metric is used (4.1).
  \item \textbf{Both of the following two bullets rest on a permutation null: the two sets are pooled, randomly re-split into groups of the original sizes, and MMD² recomputed many times, so that the spread of those values shows what MMD² looks like when there is genuinely no difference between the sets.} The standard deviation of that spread, written σ₀, is the size of a typical chance fluctuation for a given kernel, and is what both the standardization and the significance test are measured against.
  \item \textbf{Noise-floor standardization: each kernel's MMD² is divided by its own σ₀ before the kernels are summed.} Different bandwidths produce values on different numerical scales, so a plain sum lets coarse kernels swamp fine ones for reasons of scale rather than detected signal; dividing by σ₀ converts each term into a count of chance-fluctuation widths, which puts every kernel on equal footing in the sum.
  \item \textbf{Significance: the same permutation procedure is applied to the combined statistic to obtain a p-value per condition.} The p-value is the fraction of random re-splits producing a combined statistic at least as large as the observed one, reported separately for each chopper frequency and light level.
  \item \textbf{UNRESOLVED.} Both bullets above are inherited from an earlier draft and need decisions before they can be written up. First, reporting p-values is a hypothesis test, which sits awkwardly with the removal of the hypothesis-testing framing from the Introduction and Background; either the p-values stay and that framing is restored in some corrected form, or the significance bullet goes. Second, σ₀ is estimated from the data being compared, so re-estimating it per condition reintroduces exactly the data-dependence the median heuristic was rejected for, while freezing it once preserves comparability at the cost of exactness in the p-values. Third, the small-kernel sign problem noted above is a question of direction rather than magnitude, and dividing by σ₀ rescales a wrongly-signed contribution without correcting it.
\end{itemize}

\subsection{SWD}

\begin{itemize}
  \item \textbf{Wasserstein distance (WD) measures the difference between two distributions as the minimum cost of transporting the mass of one onto the other.} Picture each point set as a pile of sand: the distance is the least total work needed to reshape the first pile into the second, where work counts both how much mass is moved and how far it travels. This makes it sensitive to how far apart the two distributions are geometrically, not only to whether they overlap, which is a different kind of sensitivity from the kernel similarity used by MMD.
  \item \textbf{Computing WD directly in three dimensions is expensive, because it requires solving an optimal-transport problem between the two point sets.} In one dimension, however, the problem has a closed form: the optimal transport is simply to sort both sets and match them in order, so the distance can be read off from the sorted values in O(N log N) time. The sliced variant exploits this by reducing the 3D problem to many 1D ones.
  \item \textbf{Sliced Wasserstein Distance (SWD) projects both point sets onto a random direction, computes the 1D Wasserstein distance between the projections, and averages over many directions.} Each projection collapses the 3D points onto a line, giving two sets of scalars whose Wasserstein distance is cheap to compute; averaging over enough directions recovers a meaningful distance between the original 3D distributions:
  \item \textbf{Every symbol above is defined as follows.} X and Y are the real and synthetic point sets from 3.2, each point a 3D coordinate (x, y, s·t); θ\_ℓ is a unit vector giving the ℓ-th projection direction; θ\_ℓᵀX denotes the set of scalars obtained by projecting every point of X onto that direction; L is the number of projection directions used; W\_p is the 1D Wasserstein distance of order p between the two projected sets, computed by sorting and matching in order; and p is the order of the distance, with p = 2 the usual choice.
  \item \textbf{The projection directions are drawn by Monte Carlo sampling, which makes the reported SWD an estimate rather than an exact quantity.} Because L directions are sampled at random from the sphere rather than covering it exhaustively, repeating the computation on the same data returns slightly different values, and the size of that spread is set by L.
  \item \textbf{RESERVED: sampling-error analysis.} State how L was chosen, quantify the residual spread of SWD across repeated draws at that L, and show that the spread is small relative to the differences between conditions being reported, so that a change in SWD across conditions can be distinguished from projection noise.
\end{itemize}

\subsection{Chamfer baseline}

\begin{itemize}
  \item \textbf{Definition and role.} Symmetric nearest-neighbor distance on the same representation; the comparison point that shows what a naive point-set distance misses (distributional structure, noise sensitivity).
\end{itemize}

\section{Evaluation}
